\section{Relative boundary laws along a periodic itinerary}
\label{sec:v64-periodic-itinerary}

The relative normalization is not specific to a normal two-periodic
orbit.  We prove it for a selected nongrazing periodic orbit of a smooth
planar dispersing billiard.  The two ends now carry different half-line
actions.  There is also a linear endpoint contribution to the physical
length; it must be restored before a probability law is formed.
The subsequent principal argument treats all marked collision phases.
Alternating two-contact specializations and the analytic and lattice
extensions are retained separately in the full technical manuscript.

\subsection{The geometric Jacobi operator}

Let a periodic orbit have a fixed itinerary, with lifted reflection points
$q_i$, flight lengths $L_i>0$, curvatures $\kappa_i>0$, and outgoing
normal cosines $\nu_i>0$.  The indices are periodic modulo an even
integer $P$: an odd primitive period is doubled so that the signed
coordinate frames also return.  A lifted period is allowed to end at a
lattice translate of its initial point.  Put $L=\sum_{i=0}^{P-1}L_i$.
Every open flight is clear of unintended obstacles.  We work with a
compact smooth family of such marked orbits, with fixed $P$ and positive
lower bounds for clearance, $L_i$, $\kappa_i$ and $\nu_i$, and finite
upper bounds for the flight lengths and each separately used smooth norm.
Clearance here refers to portions outside small fixed endpoint
neighborhoods; nongrazing incidence controls those neighborhoods.
Family derivatives keep the scaled endpoint coordinates and the excess
$d=t-nL$ fixed; the physical observation time moves with the reference
period length.

Choose unit-speed signed boundary charts $r_i(s)$ centered at $q_i$.
Their tangent orientations alternate relative to the positively oriented
obstacle tangents.  Write
\[
 x_i=\nu_i s_i,\qquad
 \ell_i(x,y)=|r_{i+1}(y/\nu_{i+1})-r_i(x/\nu_i)|,
 \qquad k_i=L_i^{-1},\qquad m_i=2\kappa_i/\nu_i.
\]
The endpoint arguments in $r_i$ are signed arclengths.  Reflection gives
numbers $p_i$ such that
\[
 \ell_{i,1}(0,0)=-p_i,\qquad \ell_{i-1,2}(0,0)=p_i.
\]
Introduce the normalized edge action
\begin{equation}\label{eq:v64-edge-gauge}
 h_i(x,y)=\ell_i(x,y)-L_i+p_i x-p_{i+1}y.
\end{equation}
Thus $h_i(0,0)=0$ and $Dh_i(0,0)=0$.  The linear terms telescope;
this normalization does not change an interior reflection equation or
a mixed endpoint derivative.

\begin{lemma}[Coercivity and persistence]\label{lem:v64-geometric-jacobi}
The Hessian of an edge at the reference orbit and the interior Jacobi
operator are respectively
\begin{equation}\label{eq:v64-geometric-hessian}
 D^2h_i(0,0)=
 \begin{pmatrix} k_i+m_i/2&-k_i\\-k_i&k_i+m_{i+1}/2\end{pmatrix},
 \qquad
 (H^0x)_i=(k_{i-1}+k_i+m_i)x_i-k_{i-1}x_{i-1}-k_ix_{i+1}.
\end{equation}
They are uniformly positive.  On a sufficiently small common coordinate
box the interior Hessians are symmetric, tridiagonal, and strictly
diagonally dominant with a uniform positive margin.  The marked orbit
and these properties persist under sufficiently small smooth perturbations
of the reflecting arcs and of the lattice realization, with the lifted
itinerary fixed.
\end{lemma}
\begin{proof}
For a chord with unit direction $e$ and endpoint tangent $t$, the
same-end second derivative in signed arclength is
$(1-(e\cdot t)^2)/L$ plus the curvature contribution.  At an outgoing
reflection $e\cdot n=\nu$, with $n$ the outward obstacle normal and
$r''=-\kappa n$, that contribution is $\kappa\nu$; the incoming end
has the same sign.  The mixed derivative is
$-[t_i\cdot t_{i+1}-(e\cdot t_i)(e\cdot t_{i+1})]/L$.
The alternating tangent orientations make the two projections onto
$e^\perp$ have the same sign, so it is
$-\nu_i\nu_{i+1}/L_i$.  Dividing the endpoint derivatives by the
respective normal cosines proves the edge formula.  Summing adjacent
edge derivatives proves the second formula.

The edge quadratic form is
$k_i(x-y)^2+(m_i/2)x^2+(m_{i+1}/2)y^2$.
The finite and half-line interior forms are therefore bounded below by
$m_*\sum x_i^2$, where $m_*:=\min_i m_i>0$.
Continuity, finite periodicity, and compactness preserve a positive
row-dominance margin on one smaller box.  For persistence, apply the
finite-dimensional implicit function theorem to the reflection equations
with the prescribed periodic identification of the endpoints.  The
periodic Hessian has quadratic form
$\sum k_i(x_{i+1}-x_i)^2+\sum m_i x_i^2$, hence is invertible.
The continued orbit remains in its charts; the clearance and incidence
margins keep the chords physical.  This argument also covers a repeated
period and the prescribed deck translation.
\end{proof}

\subsection{The exact reference twist}

Set
\begin{equation}\label{eq:v64-transfer}
 Q_i=\begin{pmatrix}1+m_i/k_i&1/k_i\\m_i&1\end{pmatrix},\qquad
 \mathcal M_b=Q_{b+P-1}\cdots Q_b,
 \qquad \operatorname{tr}\mathcal M_b=2\cosh\chi.
\end{equation}
Each matrix $Q_i$ has determinant one and positive entries.  The products
for different starting phases are conjugate.  Since
$(\mathcal M_b)_{12}(\mathcal M_b)_{21}>0$ and
$\det\mathcal M_b=1$, their trace is greater than two.  In particular
$\chi>0$, and $\chi$ has a positive lower bound on the compact class.

For $N=nP$, let $W_{N,b}(u,v)$ be the stationary length along $N$
flights, starting at phase $b$, with scaled endpoint coordinates $u,v$.
Its construction on a fixed box is given below.  Define
\begin{equation}\label{eq:v64-reference-twist}
 D_{N,b}=-W_{N,b,uv}(0,0)
       =\frac{\sinh\chi}{(\mathcal M_b)_{12}\sinh(n\chi)}.
\end{equation}
In particular $D_{N,b}>0$, and its decay is comparable to $e^{-n\chi}$.
This formula is exact, rather than an unspecified exponential equivalent.

\begin{lemma}[Transfer and cofactor normalizations]
\label{lem:v64-reference}
Formula~\eqref{eq:v64-reference-twist} holds.  More generally, if $y$ is
the stationary segment and $H_N(y)$ its interior Hessian, then
\begin{equation}\label{eq:v64-cofactor}
 -W_{N,b,uv}(u,v)=
 \frac{\prod_{i=0}^{N-1}[-\ell_{b+i,12}(y_i,y_{i+1})]}{\det H_N(y)}.
\end{equation}
\end{lemma}
\begin{proof}
For the reference recurrence write
$z_i=k_{i-1}(x_i-x_{i-1})$.  Then
$z_{i+1}=z_i+m_ix_i$ and $x_{i+1}=x_i+z_{i+1}/k_i$,
which is precisely the transfer $Q_i$.  At the initial endpoint $z_0$
can be defined as $k_b(x_1-u)-m_bu$; no exterior reflection equation
is being imposed.  The left boundary derivative of the quadratic
stationary action is $-z_0-m_bu/2$.  The relation
$v=(\mathcal M_b^n)_{11}u+(\mathcal M_b^n)_{12}z_0$
therefore gives $-W_{uv}(0,0)=1/(\mathcal M_b^n)_{12}$.
Cayley--Hamilton, or diagonalization at the distinct eigenvalues
$e^{\chi}$ and $e^{-\chi}$, gives
\[
 (\mathcal M_b^n)_{12}
       =(\mathcal M_b)_{12}\frac{\sinh(n\chi)}{\sinh\chi}.
\]
The first assertion follows.  Eliminating the interior variables from
the tridiagonal Hessian gives the mixed Schur-complement entry.
The corner cofactor is the product of the adjacent mixed entries,
with the sign in \eqref{eq:v64-cofactor}.  This proves the second
assertion, including its normalization at zero.
\end{proof}

\subsection{A relative theorem on a fixed collar}

Write
\[
 E_{N,b}(u,v)=W_{N,b}(u,v)-nL+p_bu-p_bv,
 \qquad \beta_{N,b}(u,v)=\frac{-W_{N,b,uv}(u,v)}{D_{N,b}}.
\]
The superscripts $-$ and $+$ below distinguish the left future tail
from the right past tail.  They do not indicate signs of the functions.

\begin{theorem}[Periodic-itinerary relative boundary law]
\label{thm:v64-periodic-relative}
For the marked periodic orbits above there is a coordinate interval,
independent of $n$, on which the finite stationary segments exist
uniquely within the prescribed small itinerary collars.  There are
smooth strictly convex functions $S_b^-,S_b^+$ and positive smooth
functions $B_b^-,B_b^+$, with
\[
 S_b^\pm(0)=(S_b^\pm)'(0)=0,\qquad B_b^\pm(0)=1,
\]
such that, for one $\tau\in(0,1)$ and each fixed $k$,
\begin{equation}\label{eq:v64-main-relative}
 \|E_{N,b}-S_b^-\oplus S_b^+\|_{C^k}
 +\|\beta_{N,b}-B_b^-\otimes B_b^+\|_{C^k}
                  \le C_k\tau^N,\qquad N=nP.
\end{equation}
The norms include each stipulated fixed number of smooth family
derivatives.  Constants may depend on the corresponding higher norms;
no order-uniform smooth estimate is asserted.  The action error has
zero value and endpoint differential at zero and is
$O(\tau^N(|u|+|v|)^2)$.  The amplitude error vanishes at zero.
For analytic families all four limits are analytic on smaller common
complex collars.
\end{theorem}

\begin{proof}
We give the Green estimates, the construction, and the relative
normalization explicitly.  None of them uses a two-dimensional
normalizing chart.

Put $a_i=k_{i-1}+k_i+m_i$, and split the reference interior operator as
$H^0=A(I-T)$, where $A=\operatorname{diag}(a_i)$ and the two nonzero
entries in row $i$ of $T$ are $k_{i-1}/a_i$ and $k_i/a_i$.
On a finite interval or a half-line,
\[
 \|T\|_\infty\le q<1,\qquad
 (H^0)^{-1}=\sum_{r\ge0}T^r A^{-1},\qquad
 |G_{ij}|\le Cq^{|i-j|}.
\]
The last estimate holds because a nearest-neighbor walk from $i$ to $j$
has length at least $|i-j|$.  Symmetry also gives uniform $\ell^1$ bounds.
The same series is bounded on the exponentially weighted sup norm for
any weight $\rho\in(q,1)$, by summing
$q^{|i-j|}\rho^{j-i}$.  For a finite segment one uses the sum of the two
endpoint weights $\rho^i+\rho^{N-i}$.

These estimates include fixed derivatives in the marked family.
A derivative of order $r$ of a product representing a walk of length
$h$ introduces at most a fixed polynomial in $h$, times a bound
$C_rq^{h-r}$.  Alternatively one may differentiate the resolvent identity.
A strict enlargement of the exponential base absorbs that polynomial.
The finite Green kernel and the left half-line kernel also satisfy
\begin{equation}\label{eq:v64-green-reflection}
 |G^N_{ij}-G^-_{ij}|\le Cq^{2N-i-j}\quad (i,j<N),
 \qquad |G^N_{ij}|\le Cq^{|i-j|}.
\end{equation}
Indeed a walk contributing to their difference must reach the removed
right boundary and return.  Its length is at least $2N-i-j$.
The reflected right version and the fixed differentiated versions follow
in the same way, with a slower base when necessary.  The diagonal
coefficients on retained sites are identical; imposing a Dirichlet
boundary deletes walks, not those diagonal coefficients.

Solve the finite reflection equations as their reference linear
solution minus $G^N$ times the nonlinear gradient remainder.  That
remainder is local, vanishes quadratically at zero, and has derivative
$O(r)$ in a box of radius $r$.  The weighted Green bounds make this a
contraction with constant at most $1/2$ on a ball of radius
$C(|u|+|v|)$ when the common endpoint interval is sufficiently small.
This constructs the physical segment with
\[
 |y_i|\le C(|u|\rho^i+|v|\rho^{N-i}).
\]
Strict convexity gives uniqueness among all segments in the small box,
not just in the contraction ball.  The same construction on the future
and past half-lines gives $x_i^-(u)$ and $x_i^+(v)$, decaying as
$C|u|\rho^i$ and $C|v|\rho^i$.  Differentiating the fixed-point
equation leaves the same invertible operator on its highest derivative;
the other terms are local products of lower derivatives.  Induction
therefore gives the weighted bounds at every fixed order.

Define the half-line actions by the absolutely convergent sums
\begin{equation}\label{eq:v64-half-actions}
 \begin{split}
 S_b^-(u)&=\sum_{i\ge0}h_{b+i}(x_i^-(u),x_{i+1}^-(u)),\\
 S_b^+(v)&=\sum_{i\ge0}h_{b-i-1}(x_{i+1}^+(v),x_i^+(v)).
 \end{split}
\end{equation}
Every summand is quadratic in its exponentially decaying arguments.
The fixed differentiated sums converge as well.  In the first variation
of a finite truncation the interior terms cancel.  The remaining remote
boundary term tends to zero by the weighted estimates, so these sums
have the stated endpoint derivatives.  The positive edge forms in
Lemma~\ref{lem:v64-geometric-jacobi} bound each minimized half-line
energy below by a fixed positive multiple of its boundary value squared.
Their Hessians at zero are consequently positive.  Shrinking the interval
proves strict convexity and the stated anchoring.

For the action comparison, glue the two half-lines by
$z_i=x_i^-(u)+x_{N-i}^+(v)$ and overwrite the two endpoint values by
$u,v$.  Locality and $Dh_i(0,0)=0$ imply that the interior residual
has $\ell^1$ norm at most $C(1+N)\rho^N$: every interaction contains
one factor from each end.  The endpoint overwrites have the same order.
The averaged Hessian between $z$ and the finite stationary segment has
uniformly bounded $\ell^1$ inverse by diagonal dominance and symmetry.
Hence $\|y-z\|_1\le C(1+N)\rho^N$.  Subtracting and differentiating
these equations proves the same assertion at each fixed order, with
only a polynomial loss in $N$.  Summing edge interactions and the
omitted half-line tails now proves the action part of
\eqref{eq:v64-main-relative}.

For the amplitude, let $G_N=(H_N(0))^{-1}$ and
$\Delta H_N=H_N(y)-H_N(0)$.  The exact cofactor identity gives
\begin{equation}\label{eq:v64-relative-log}
 \log\beta_{N,b}=
 \sum_{i=0}^{N-1}\log\frac{-\ell_{b+i,12}(y_i,y_{i+1})}{k_{b+i}}
            -\log\det(I+G_N\Delta H_N).
\end{equation}
In particular the exponentially small $D_{N,b}$ has already been
cancelled exactly.  Tridiagonality and endpoint localization give
$\|\Delta H_N\|_1\le C(|u|+|v|)$, where the norm is the trace norm:
the sum of absolute matrix entries is an upper bound for that norm.
All its fixed derivatives have uniformly bounded trace norms.
Take the endpoint interval small enough that
$\|G_N\Delta H_N\|<1/2$, uniformly in $N$.

On the future half-line define $\Delta H_b^-(u)$ by the same rule,
and put
\begin{equation}\label{eq:v64-half-amplitude}
 \begin{split}
 \log B_b^-(u)
  ={}&\sum_{i\ge0}\log
       \frac{-\ell_{b+i,12}(x_i^-(u),x_{i+1}^-(u))}{k_{b+i}}\\
    &-\sum_{r\ge1}\frac{(-1)^{r+1}}r
                   \operatorname{tr}[(G_b^-\Delta H_b^-(u))^r].
 \end{split}
\end{equation}
The analogous reversed formula defines $B_b^+$.  Each edge sum is
absolutely convergent.  The perturbation is trace class and the operator
norm is less than $1/2$, so the trace series and its fixed derivatives
converge.  Its exponential is positive for real arguments, and equals
one at zero.  This specifies the determinant normalization completely.

Here is the two-end comparison for \eqref{eq:v64-relative-log}.
For $N\ge6$ retain the first and last $r=\lfloor N/3\rfloor$
interior sites, and delete all other perturbation entries.  The
trace-norm cost is $C_k(1+N)^{a_k}\rho^r$ after $k$ fixed
derivatives.  On either retained block the gluing estimate replaces
the finite segment by its own half-line with exponentially small
entrywise sum.  The identity $\det(I+AB)=\det(I+BA)$ reduces a
finitely supported perturbation to those two blocks.
Equation~\eqref{eq:v64-green-reflection} compares the diagonal Green
compressions to their half-line compressions; the opposite-end Green
blocks are $O(q^{N-2r})$, up to a polynomial factor in $N$.
Multiplication by the bounded trace-norm perturbation preserves these
bounds.  Extending each retained block to its half-line costs
$O(\rho^r)$ again.

For trace-class $T,\widetilde T$ with norms at most $q_1<1$, the
exact telescoping expansion of their powers gives
\[
 |\operatorname{tr}(T^h-\widetilde T^h)|
       \le hq_1^{h-1}\|T-\widetilde T\|_1.
\]
Summing with coefficient $1/h$ proves the required continuity of the
logarithmic determinant.  At each fixed differentiated order the series
has an additional fixed polynomial in $h$, still summable.  Thus the
finite logarithm tends relatively to the sum of the two logarithms in
\eqref{eq:v64-half-amplitude}.  The edge-product comparison follows
from the same gluing and tail bounds.  Exponentiation proves the
amplitude part of \eqref{eq:v64-main-relative}.

Choose one $\tau<1$ slower than the Green, gluing, and block-truncation
rates.  The strict margins absorb every separately fixed polynomial
loss; the constants, not this choice of $\tau$, depend on the order.
Finitely many shorter segments are absorbed by enlarging those constants.
At zero the actions have zero value and differential, and the normalized
twists and amplitudes equal one.  Taylor's integral remainder gives the
additional quadratic vanishing.  Finally the contraction, summable
perturbations, and logarithmic series extend to uniform smaller complex
boxes for analytic families.  Their locally uniform convergence proves
the analytic assertion.
\end{proof}

\subsection{Nonlinear transmission and the stable return}

Define the connected endpoint action and the two increasing coordinates
by
\[
 \begin{split}
 C_{N,b}(u,v)&=W_{N,b}(u,v)-W_{N,b}(u,0)
                         -W_{N,b}(0,v)+W_{N,b}(0,0),\\
 \zeta_b^\pm(x)&=\int_0^x B_b^\pm(t)\,\dd t.
 \end{split}
\]

\begin{corollary}[Relative nonlinear transmission]
\label{cor:v64-transmission}
On a smaller fixed interval,
\begin{equation}\label{eq:v64-transmission}
 -\frac{C_{N,b}(u,v)}{D_{N,b}}
       =\zeta_b^-(u)\zeta_b^+(v)+R_{N,b}(u,v),\qquad
 \|R_{N,b}\|_{C^k}\le C_k\tau^N,
\end{equation}
with $|R_{N,b}(u,v)|\le C\tau^N|uv|$.
Let $T_b^-$ be the next-period position on the future stable half-line,
and $T_b^+$ the previous-period position on the past half-line.  Then
\begin{equation}\label{eq:v64-period-linearizer}
 \zeta_b^\pm(T_b^\pm(x))=e^{-\chi}\zeta_b^\pm(x).
\end{equation}
Each $\zeta_b^\pm$ is the unique differentiable local linearizing
coordinate with value zero and derivative one at zero.
\end{corollary}
\begin{proof}
Twice integrating the mixed derivative, with oriented integrals on
negative intervals, gives the exact identity
\[
 -C_{N,b}(u,v)/D_{N,b}
             =\int_0^u\int_0^v\beta_{N,b}(s,t)\,\dd t\,\dd s.
\]
The relative theorem proves \eqref{eq:v64-transmission}, including
its endpoint vanishing and fixed mixed derivatives.  Notice that the
unnormalized remainder has the smaller size
$O(D_{N,b}\tau^N|uv|)$; an absolute action estimate alone does not
give this transmission coefficient.

For the future return concatenate a $P$-flight segment with an
$N$-flight tail.  The joining endpoint $w_N(u,v)$ is uniquely determined
by stationarity and tends to $T_b^-(u)$ with every fixed derivative.
The denominator in its implicit equation is the sum of two positive
same-end Hessians.  Exact stationary differentiation gives
\[
 -W_{N+P,b,uv}(u,v)
       =[-W_{N,b,12}(w_N(u,v),v)]\,\partial_u w_N(u,v).
\]
Since $D_{N+P,b}/D_{N,b}\to e^{-\chi}$, relative factorization
and cancellation of the positive right factor yield
\[
 B_b^-(T_b^-(u))(T_b^-)'(u)=e^{-\chi}B_b^-(u).
\]
Integration gives the minus identity in \eqref{eq:v64-period-linearizer}.
Reversing the itinerary proves the plus identity.  In particular both
return maps have derivative $e^{-\chi}$ at zero and are contractions
on a smaller interval.  If $\widetilde\zeta$ is another normalized
differentiable linearizer, then
$h=\widetilde\zeta\circ(\zeta_b^\pm)^{-1}$ satisfies
$h(\lambda x)=\lambda h(x)$, $h(0)=0$, $h'(0)=1$, where
$\lambda=e^{-\chi}$.  Thus
$h(x)=\lambda^{-n}h(\lambda^nx)\to x$, proving uniqueness.
\end{proof}

\subsection{The physical window law}

The gauge in \eqref{eq:v64-edge-gauge} is useful analytically, but it is
not a change of the clock.  Define the physical boundary actions
\begin{equation}\label{eq:v64-physical-actions}
 \Sigma_b^-(u)=S_b^-(u)-p_bu,\qquad
 \Sigma_b^+(v)=S_b^+(v)+p_bv.
\end{equation}
In general they do not have critical points at zero.  Their linear parts
are the oblique endpoint momenta of the reference orbit.

\begin{theorem}[Charged phase-volume law for a repeated itinerary]
\label{thm:v64-physical-law}
Let $A$ be the free area of the periodic billiard cell and normalize its
unit-speed phase volume by $2\pi A$.  Fix a sufficiently small positive
offset interval $[d_-,d_+]$.  There is a smaller closed endpoint interval
$I$, independent of $N=nP$, such that the following experiment is physical.
The first collision is in the phase-$b$ collar; there are the selected
$N$ successive flights and no further collision by time $t_N=nL+d$;
the two recorded endpoints lie in $I^2$.  The intervening collisions
are required to remain in the prescribed itinerary collars.

Its endpoint measure, before conditioning on success, has density
\begin{equation}\label{eq:v64-phase-density}
 \frac{D_{N,b}}{2\pi A}\,
 \beta_{N,b}(u,v)
       \{d-E_{N,b}(u,v)+p_bu-p_bv\}_+,\qquad (u,v)\in I^2.
\end{equation}
On this interval the expression inside braces is bounded below by
$d_-/2$ and above by a constant smaller than either adjacent free-flight
length.  Write $Z_{N,b,d}$ for the integral of the numerator after the
constant $D_{N,b}/(2\pi A)$ is removed.  The success probability is
exactly $D_{N,b}Z_{N,b,d}/(2\pi A)$ and is bounded above and below by
positive constant multiples of $e^{-n\chi}$.

The conditional endpoint density converges, in every fixed $C^k(I^2)$
norm including the stated family derivatives, to
\begin{equation}\label{eq:v64-limit-law}
 f_{b,d}(u,v)=Z_{b,d}^{-1}B_b^-(u)B_b^+(v)
                \{d-\Sigma_b^-(u)-\Sigma_b^+(v)\},
 \qquad \|f_{N,b,d}-f_{b,d}\|_{C^k}\le C_k\tau^N.
\end{equation}
All rejected independent phase preparations remain charged through the
success probability.  The result is a law for this selected windowed
itinerary, not an identification with a maximal-collision event.
\end{theorem}
\begin{proof}
Clearance and nongrazing incidence imply that the actual preceding and
following free flights of every sufficiently small stationary segment
have a common positive lower bound $l_*$.  Choose $d_+<l_*/4$.
The action has a uniform quadratic bound on the fixed box, whereas
$|p_b|$ is uniformly bounded.  Shrink $I$ until
$|E_{N,b}-p_bu+p_bv|\le d_-/2$, uniformly in $N$.
Then the available residual time
$r=t_N-W_{N,b}=d-E_{N,b}+p_bu-p_bv$ lies in
$[d_-/2,3d_+/2]$, below $l_*$.  Any split $r=s+(r-s)$,
$0\le s\le r$, lies strictly within the adjacent flights (apart from
measure-zero collision endpoints).  It gives exactly the required first
and last collision times.  Conversely a successful trajectory has this
unique first collision, its unique small stationary bridge, and such a
split.  Thus no multiplicity is introduced by repeating the word.

If $x$ is scaled boundary arclength, its canonical tangential momentum
is the ordinary arclength momentum divided by $\nu_b$.  Hence
$dx\wedge d\pi$ is the usual billiard flux form.  The endpoint
generating action gives $\pi=-W_{N,b,u}$, so that the absolute
Jacobian to $(u,v)$ is $-W_{N,b,uv}>0$.  The full phase form is this
flux times $ds$.  Integrating over $0\le s\le r$ proves
\eqref{eq:v64-phase-density}, with precisely the normalization
$2\pi A$.  In particular the endpoint gauge terms have been restored.

The residual interval has a positive floor.  The finite and limiting
amplitudes have positive upper and lower bounds on a common box, by
\eqref{eq:v64-relative-log} and \eqref{eq:v64-half-amplitude}.
Consequently $Z_{N,b,d}$ and $Z_{b,d}$ have positive uniform floors
and finite uniform upper bounds.  The relative theorem compares their
unnormalized integrands and fixed derivatives at rate $C_k\tau^N$.
Integration on the fixed square and division by these floors prove
\eqref{eq:v64-limit-law}.  Finally
\eqref{eq:v64-reference-twist} gives the stated charged probability.
At no stage is an absolute action error divided by that probability.
\end{proof}

\subsection{A nonnormal orbit and the alternating specialization}

\begin{proposition}[A three-obstacle realization]
\label{prop:v64-triangle}
There is a periodic dispersing table with a nongrazing three-periodic
orbit to which Theorems~\ref{thm:v64-periodic-relative}
and~\ref{thm:v64-physical-law} apply, and an open family of such
realizations with unequal flight lengths and incidence angles.
\end{proposition}
\begin{proof}
In a square cell of side $12$ place three unit disks whose centers form
an equilateral triangle of side $D=5/2$, centered in the cell.
The disks and their translates are disjoint.  At each disk choose the
point on its radius directed toward the triangle center.  These three
points form an equilateral triangle of side
$\ell=D-\sqrt3>0$.  Its inscribed three-flight path satisfies reflection:
the outward disk normal bisects the two outgoing directions, with normal
cosine $\nu=\sqrt3/2$.  The distance from the opposite disk center to
a whole side line is $\sqrt3 D/2-1/2>1$; hence no third disk blocks
that side.  Neighboring cells are much farther away.  The flights are
therefore clear and the incidence is not normal or grazing.

Take the oriented period $P=6$.  Then $k_i=1/\ell$,
$m_i=2/\nu$, and
\[
 \chi=6\operatorname{arcosh}(1+\ell/\nu).
\]
For explicit nonsymmetric realizations, choose radii $R_0,R_1,R_2$ close to one and prescribe a nearby scalene
contact triangle.  Place center $i$ at distance $R_i$ behind vertex $i$
along its interior-angle bisector.  The triangle then satisfies reflection
by construction, has unequal sides and, choosing its three angles
distinct, unequal normal cosines.  Disjointness and clearance persist by
the strict margins just checked.  This gives the asserted open family
by the same persistence lemma.
\end{proof}

For a normal alternating channel take $P=2$, $L_0=L_1=g$ and
$\nu_0=\nu_1=1$.  Then $p_b=0$, $\chi=2\gamma$, and
\eqref{eq:v64-reference-twist} reduces to
$D_{2n,b}=a_b/\sinh(2n\gamma)$.  Reversal identifies the two
half-line actions and amplitudes.  Passing from arclength to the fixed
transverse graph chart changes an amplitude by the corresponding
endpoint Jacobian, normalized by its value at zero.  Thus the relative
law specializes to Theorem~\ref{thm:v4-factorization}, not to a different
normalization of it.  The intrinsic scalar coordinate in
Corollary~\ref{cor:v64-transmission} specializes to
Theorem~\ref{thm:v14-intrinsic-density}.

The extension above concerns the geometric forward mechanism and its
physical normalization.  It neither assumes nor proves that the two
boundary actions of a general itinerary can be inverted by the
alternating two-contact block formula.  The next sections establish the general periodic finite-jet and smooth
inverses in measured graph coordinates.  The alternating inverse, analytic
extensions and global lattice reconstruction in the full technical
manuscript retain their separate hypotheses.  The finite-preparation
consequences of the smooth inverse are proved in the principal article
and its two-offset appendix.

\input{article/10k_relative_decoupling_v76}
